
\documentclass[a4paper,12pt]{article}

\input{glyphtounicode}
\pdfgentounicode=1

\usepackage{mathptmx}
\usepackage{fancyhdr}
\usepackage{hyperref}
\usepackage[english]{babel}
\usepackage[utf8]{inputenc}
\usepackage[version=4]{mhchem}
\usepackage[usenames,dvipsnames]{color}
\usepackage[
  a4paper,
  top    = 2.00cm,
  bottom = 2.00cm,
  left   = 1.70cm,
  right  = 1.70cm
]{geometry}

\renewcommand{\baselinestretch}{1.5}
\renewcommand{\headrulewidth}{0pt}
\renewcommand{\footrulewidth}{0pt}

\hypersetup{
    colorlinks=true,
    linkcolor=RoyalBlue,
    urlcolor=RoyalBlue,
    citecolor=RoyalBlue,
    anchorcolor=RoyalBlue
}

\urlstyle{same}
\pagestyle{fancy}
\fancyhf{}
% \fancyfoot[R]{\thepage}

\begin{document}

\begin{center}
    {\Huge \scshape {\fontsize{25}{30}\selectfont{Reza}} {\fontsize{25}{30}\selectfont{\textbf{Aliasgari Renani}}}} \\ \vspace{3pt}
    {\small {07/12/2025}} $|$
    {\small {Research Statement}} $|$
    {\small {Hamburg U., Heidelberg U., Munich U., RWTH Aachen U.}} $|$
    {\small {PhD Position}} \\
    {\small {Foundation Models in Astrophysics and Particle Physics}}
    \vspace{-10pt}
\end{center}
\vspace{-10pt}
\noindent\rule{\textwidth}{0.5pt}

My research objective is to contribute to the advancement of particle physics through development of foundation models, in particular by applying knowledge of radiation effects on hardware and creating FPGA based systems for efficient data processing in machine learning applications. The PhD position within the SciFM project in Foundation Models in Astrophysics and Particle Physics at RWTH Aachen University offers an excellent opportunity for me to pursue this goal. I am now completing the final year of my Masters degree in Plasma Physics at Moscow Institute of Physics and Technology, where my research on semiconductor devices exposed to different kinds of radiation and FPGA development has strengthened my interest in studying particle physics with machine learning techniques. I am especially motivated by the possibility of contributing to interpretable models within a multidisciplinary collaboration and using hardware expertise for scientific discoveries.

\vspace{5pt}

I joined the laboratory of Dr. Koveshnikov in the Russian Academy of Sciences over two years ago to study radiation effects on microelectronic structures, focusing on electrical characterization of semiconductors and \ce{SiO2} based MOS devices. We studied the impact of electron and gamma irradiation and hydrogen plasma treatment on MOS devices, identified traps by location (oxide, semiconductor, interface) and carrier type (electrons, holes). This work has resulted in a publication as first author\footnote{R. Aliasgari Renani, O.A. Soltanovich, M.A. Knyazev, S.V. Koveshnikov, Investigation of low energy electron irradiated \ce{SiO2} based MOS devices by C-V and TSC techniques, Russian Microelectronics, 2023} and provides insight into device stability under bias stress, resistance to radiation, and density of oxide traps at the interface between dielectric and semiconductor. I also developed optimized MATLAB applications to automate the experimental characterization techniques, which reduced run times and created measurement pipelines.

\vspace{5pt}

After my undergraduate degree, I joined the System on Chip Development Laboratory and the Laboratory of Plasma Systems under the supervision of Dr. Vasiliev to develop FPGA based systems and to study radiation and plasma effects on these structures. I implemented video processing algorithms by manual Verilog coding and by generating HDL from Simulink models. I verified functionality of the algorithms using testbenches and Python automation scripts. In the Plasma Lab, we ran experiments where an FPGA with an attached camera and a synthesized image processing design was placed in a plasma system from electron beam with oxygen at low pressure and exposed to electron irradiation and X rays generated from a tungsten target. The FPGA output signal was monitored while thermal cycling was applied and preliminary functionality tests were performed.

\vspace{5pt}

My experience in these projects has equipped me with a strong foundation in handling complex data from radiation experiments and in designing hardware for robust performance under stress. These skills align well with the demands of developing foundation models in particle physics, where large datasets from simulations and experiments require efficient processing, automation of analysis pipelines, and consideration of hardware reliability in high radiation settings typical of accelerators.

\vspace{5pt}

My technical background maps directly to the methods and requirements of the SciFM PhD position at RWTH Aachen University, which emphasizes programming and machine learning techniques for building foundation models in particle physics. In programming, I possess advanced proficiency in MATLAB and Simulink for automation of experimental techniques and data processing, as demonstrated by my development of applications that control instruments like Keithley SourceMeter and Zurich Instruments MFIA to conduct measurements such as thermally stimulated current and capacitance voltage. This expertise extends to Python, where I have created scripts for simulation automation, data analysis, and verification of digital designs, skills essential for implementing and training machine learning models on large datasets. My work with NumPy, SciPy, and OpenCV for image processing algorithms, including color segmentation and edge detection, provides a basis for handling multidimensional data similar to particle collision events, enabling me to contribute to model architectures that process detector outputs. In hardware design, my Verilog and RTL skills for FPGA development, including synthesis with Vivado and implementation of fixed point computations, align with the need for efficient inference in foundation models, particularly for deployment in resource constrained environments. Additionally, my knowledge of radiation effects on semiconductors and plasma interactions informs the creation of robust models that account for real world detector noise and degradation in high energy physics experiments. I am prepared to collaborate on interpretable and open access models within the multidisciplinary SciFM network.

\vspace{5pt}

In alignment with the SciFM focus on developing large scale foundation models for particle physics at RWTH Aachen University, I propose to pursue research that integrates my hardware and data processing expertise with machine learning to create powerful, interpretable models for analyzing collision data. Building on my FPGA background, a key goal is to optimize these models for hardware acceleration, implementing efficient inference pipelines on FPGA platforms to handle real time processing of high volume data from accelerators, while accounting for radiation induced effects to enhance model robustness in experimental settings. I intend to validate these models against real data from collaborations, collaborating with the SciFM network to bridge particle physics with astrophysics through shared machine learning approaches. Under the supervision of Prof. Michael Krämer, this work would contribute to open access models that advance precision measurements and searches for physics beyond the standard model.

\vspace{5pt}

My experience with irradiated semiconductors, plasma systems and FPGA design motivated me to continue my PhD studies at RWTH Aachen University. I aim to engage in development of large scale foundation models for particle physics to address fundamental questions. My background will enable me to contribute to the projects on machine learning models that are powerful and interpretable, by applying my knowledge of radiation effects and FPGA implementation for efficient model deployment in high energy environments. I am enthusiastic to work under the guidance of Prof. Michael Krämer at RWTH Aachen University to integrate machine learning techniques with hardware solutions for advancing particle physics research.

\vspace{5pt}

Having been an international student and having collaborated across multiple scientific institutions, I am comfortable working between partner laboratories and I would welcome the opportunity to collaborate across sites in Germany to advance this research. I am driven by research that couples numerical methods with experiments, and I look forward to leveraging my skills in FPGA development for efficient model inference and to validate the models against data from particle physics experiments, which would ultimately contribute to breakthroughs in understanding fundamental physics. It would be an honor to study and work at RWTH Aachen University in Germany, a country with longstanding traditions of scientific and cultural excellence.

\end{document}